\documentclass[11pt,letterpaper]{article}
\usepackage[T1]{fontenc}
\usepackage{lmodern,amsmath,amssymb,booktabs,longtable,graphicx}
\usepackage[margin=1in]{geometry}
\usepackage[hidelinks]{hyperref}
\newtheorem{theorem}{Theorem}
\newtheorem{proposition}[theorem]{Proposition}
\newtheorem{lemma}[theorem]{Lemma}
\newtheorem{corollary}[theorem]{Corollary}
\newtheorem{remark}[theorem]{Remark}
\renewcommand{\square}{\vcenter{\hbox{\rule{.72ex}{.72ex}}}}
\newenvironment{proof}{\par\smallskip\noindent\textit{Proof.}\ }{\nobreak\hfill\rule{.72ex}{.72ex}\par\smallskip}
\newcommand{\Z}{\mathbb Z}
\newcommand{\N}{\mathbb N}
\newcommand{\authorline}{Iaroslav Argunov}
\newcommand{\paperhead}[1]{\title{#1}\author{\authorline}\date{18 November 2026}\maketitle}
\renewcommand{\arraystretch}{1.12}
\setlength{\parindent}{1.1em}
\raggedbottom

\hypersetup{pdfauthor={Iaroslav Argunov},pdftitle={Congruences for closed meanders from rotations and curve interchange}}
\begin{document}
\paperhead{Congruences for closed meanders from rotations and curve interchange}
\begin{abstract}
Let \(M_n\) be the number of ordered pairs of noncrossing matchings on \(2n\) cyclically labelled sites whose union is connected. Two finite group actions determine a residue modulo \(8n\) at odd orders where the smaller fixed sets can be counted. In particular, \(M_{39}\equiv16\pmod{312}\). We give two ways to determine the projective term modulo eight, one by enumerating self-dual diagrams. The latter also characterizes exactly the odd orders admitting a projective word whose square reverses both curves.
\end{abstract}
Draw the two matchings in opposite hemispheres. Their common equator and their alternating circuit are the two curves of an ordinary meander. Retain the equatorial labels and the order of the hemispheres. The closed sequence in this convention is A005315. La Croix describes the symmetries of the two-curve model \cite{lacroix}. Howroyd gives a relation between reflection classes and the closed count \cite{how}. We combine the action on ordinary pairs with a second action on its projective fixed set. The accompanying exact tables give the projective count \(P_n\) and the symmetric marked count \(D_n=2K_{n-1}\) for \(n\ge2\). At the separate base case \(D_1=K_0=1\).

\section{Eight symmetries of a connected pair}
Put \(H(i)=i+n\bmod 2n\). Let \(X\) exchange the two ordinary noncrossing matchings on \(2n\) labelled sites and let \(R(i)=-1-i\bmod 2n\). For odd \(n>1\), \(L=HR\) is a reflection fixing two sites. The commuting involutions \(X,H,R\) act on ordered connected pairs. At \(n=39\) their fixed sets give
\begin{center}\small
\begin{tabular}{llr}\toprule
Element & Fixed configurations & Count\\\midrule
identity & ordinary closed diagrams & \(M_{39}\)\\
\(X,L\) & common isolated chord, or identical matchings & 0 each\\
\(H\) & labelled stamp foldings & \(23\,438\,785\,044\,156\,653\,184\)\\
\(R\) & open meanders on 39 crossings & \(64\,477\,712\,119\,584\,604\)\\
\(RX\) & semi-meanders on 39 crossings & \(600\,994\,488\,311\,709\,056\)\\
\(HX\) & antipodal connected matchings & \(1\,388\,207\,765\,065\,576\,684\)\\
\(LX\) & labelled symmetric open diagrams & \(4\,790\,198\,939\,133\,878\,624\)\\
\bottomrule\end{tabular}
\end{center}
For \(H\), odd parity forces one diameter in each hemisphere. After quotienting, the two uncovered endpoints of the matchings are the ends of one path across the ordered levels. Reading the levels along this path gives a folding of a labelled strip, with the upper colour choosing the first end. For \(R\), reflect a half-diagram to complete it. A reflection of a connected alternating cycle through gaps fixes exactly two edges. Cutting there yields a single open meander on \(n\) sites; reflection is inverse to the cut. For \(RX\) the lower matching is determined by the upper one. The reflection edges form the fixed rainbow. For \(HX\), use \((AH)^2=A(HAH)\): a component of \(A\cup H\) with \(2r\) sites gives cycles of length \(r\) of \(AH\). The square gives a connected upper/lower pair exactly when there is one component and \(r=n\) is odd. For \(LX\) the reflection has two fixed vertices. Its edges and \(A\) form an alternating path and possibly separate cycles. The path uses the entire diagram precisely in the connected case. Projecting from one fixed vertex of the sphere gives the labelled centrally symmetric open meanders on \(2n-1\) crossings, whose number is \(2K_{n-1}\) \cite[Prop.~6.1]{legendre}. The folding and meander conventions are those of Sawada and Li \cite{folding}. The published values in the first three nonzero rows are given in A000136, A005316 and A060206 \cite{tables,hogan}.

The two newly evaluated entries in the table are respectively \(4\) and \(0\) modulo eight. Summing fixed counts and dividing into orbits gives \(M_{39}\equiv0\bmod8\). By the face-count formula below, for the ordinary rotation subgroups of orders 3 and 13 the fixed counts are \(640\,398\,992\,404\) and \(68\). Orbit sizes give \(M_{39}\equiv1\bmod3\) and \(3\bmod13\). The three coprime moduli therefore give
\[
 \boxed{M_{39}\equiv16\pmod{312}}.
\]
This congruence is a check on a future exact enumeration. Applying the same two group actions at other odd orders gives the following additional entries. The residue in the second column refers to the full labelled closed count. The last column follows by adding the symmetric labelled count and dividing by four in Legendre's shape formula.
\begin{center}\begin{tabular}{rrr}\toprule
Order & \(M_n\bmod 8n\)&\(a(2n-1)\bmod 2n\)\\\midrule
39 & 16 & 60\\
41 & 84 & 43\\
43 & 260 & 37\\
45 & 72 & 82\\
\bottomrule\end{tabular}\end{center}


\section{The rotational term at odd order}
Fix the equatorial gap immediately before label zero. In either disk the \(m\) matching arcs cut out \(m+1\) faces. Root the dual tree at the face touching the distinguished gap. Write \(e_A,o_A\) for the numbers of even- and odd-depth faces of a matching \(A\), and similarly for \(B\). Put
\[
  F_m=\sum_{(A,B)\text{ connected}}(e_Ao_B+o_Ae_B).
\]
For an odd divisor \(d>1\) of \(n\) a rotation of order \(d\) fixes exactly \(F_{n/d}\) ordinary connected pairs. Indeed, no diameter can occur, since a second rotated diameter would cross it. All chord orbits have length \(d\). Quotient either disk by this rotation to obtain a matching on \(2n/d\) sites with a marked face containing the rotation centre. Conversely the marked face determines, for each arc, its boundary interval away from the centre. Lifting these nested or disjoint intervals recovers the unique invariant matching on the original labels. A disconnected quotient cannot lift to a single component. A connected quotient is a simple loop on the sphere punctured at the two centres. Its lift is connected precisely when this loop separates the punctures, in which case it winds once up to sign. The two root faces touch through the distinguished gap. Crossing an arc changes side of the loop, so separation requires the marked faces to have opposite depth parities. This gives the formula. The cyclic group generated by a shift by two labels therefore gives
\[
 M_n\equiv-\sum_{1<d\mid n}\varphi(d)F_{n/d}\pmod n .
\]
Since \(n\) is odd, this combines with the residue modulo eight without ambiguity. The supplement includes two ways of computing the face sums, by forward transfer and by joining half-diagrams. Their entries agree through order 13 and at order 15, where \(F_{15}=78\,448\,367\,021\,148\).

\section{Symmetries in odd projective order}\label{app:residue}

Throughout this section \(n>1\) is odd. Write \(H(i)=i+n\) on the
\(2n\) boundary labels and let \(P_{n,h}\) denote the contribution of
height \(h\) to \(P_n\). We use \(M_j\) for the ordinary closed
meandric number. On \(2j\) labels let \(D_j\) count noncrossing
matchings \(C\) for which \(C\cup tCt\) is connected, where \(t\)
is a reflection through two opposite vertices.

\subsection*{A free action above height one}
For any division of the boundary into two half intervals,
\[
  4\mid P_{n,h}\quad(h>1),\qquad
  P_n\equiv M_{(n+1)/2}\pmod 4.                         \tag{A1}
\]
Indeed, take the halves \(0,\dots,n-1\) and \(n,\dots,2n-1\).
The half turn \(H\) and the reflection \(z(i)=n-1-i\) preserve
height and generate a group of order four. If \(HAH=A\), any arc
internal to a half belongs to an isolated four-cycle of \(A\cup H\).
If all arcs pass between the halves, their ends pair in reverse
order and the middle arc is an edge of \(H\). Neither case gives
a connected object. The two fixed sites of \(z\) must be matched
to each other in a \(z\)-invariant matching, again along an edge of
\(H\).

The remaining involution \(Hz\) exchanges the halves and reverses
their orders. In an invariant matching it fixes each of the \(h\)
through arcs setwise, since their ends pair in reverse order. It
also fixes exactly one edge of \(H\) setwise. An involution of a
cycle of length \(2n>2\) without fixed vertices fixes at most two
edges. Thus the action is free for \(h>1\). At height one, append
a site to each half and match it to the unique through end of that
half. In place of the through edge this inserts a path through the
two extra sites and their antipodal edge. After identifying paired
sites in opposite halves, the two local matching systems form an
ordinary meander of order \((n+1)/2\). Removing the extra sites
recovers the original diagram. This proves (A1).

\subsection*{Reversing and interchanging the two curves}
View an ordinary meander as a pair of simple circles on the sphere.
Keep the first circle's labels in cyclic order and record the labels
on the second as a word \(p(0),\ldots,p(2n-1)\). Start at \(p(0)=0\)
and first follow an upper arc. Then \(p\) preserves parity. The upper
and lower arcs join successive word entries, with the upper arcs
starting at even indices. In particular, the word determines both
matchings.

The projective objects correspond to the words with \(pH=Hp\).
Indeed, write \(A\) for the upper matching. The lower matching is
\(HAH\). The products for \(A\cup H\) and \(A\cup HAH\) are \(AH\)
and \((AH)^2\). Each component of a union of perfect matchings
gives two orbits of its product. When the first graph is connected,
its two orbits have odd length \(n\) and remain orbits on squaring.
Conversely, squaring cannot decrease the number of orbits.
On the connected road of \(A\cup HAH\), \(H\) is a fixed-point-free
involution that exchanges the two edge colors. It fixes no edge,
so it shifts the road by \(n\) positions. Conversely, a word
commuting with \(H\) has its upper and lower arcs interchanged by
\(H\), since their starting indices differ by the odd number \(n\).

Set \(r(i)=-i\) modulo \(2n\). Reversing the two circle orientations
gives \(p\mapsto rp\) and \(p\mapsto pr\). Interchanging the two
oriented circles gives \(p\mapsto p^{-1}\), with the first outgoing
arc specifying the upper side in the new drawing. Each operation
preserves the projective condition. In this class let
\[
 E_n=\#\{p:p=p^{-1}\},\qquad Q_n=\#\{p:p^2=r\},\qquad
 k=(n+1)/2 .
\]
The two reversals and interchange define an action of the dihedral
group of order eight on the words. The action may have a kernel.
The orbit-counting formula gives
\[
 P_n+D_k+2E_n+2Q_n\equiv0\pmod 8.                      \tag{A2}
\]

To see the \(D_k\) term, simultaneous reversal fixes precisely the
words commuting with \(r\). On such a word \(r\) interchanges
upper and lower arcs, so \(rAr=HAH\). This makes \(A\) invariant
under the gap reflection \(Hr\). After a rotation of the cut,
this is the reflection \(Hz\) in the proof of (A1). An invariant
connected \(A\) has height one in that cut. In the ordinary
order-\(k\) construction its second matching is the reflection of
its first. The reflection reverses the original \(n\) sites of a half,
and fixes both the middle site and the appended site. Hence there are
exactly \(D_k\) possibilities. Neither single reversal fixes a word.
Interchange alone and interchange followed by both reversals are
conjugate, giving \(E_n\) each. Interchange followed by either single
reversal requires \(p^2=r\), giving the two \(Q_n\) terms.

In fact \(Q_n=0\) unless \(n\equiv1\pmod8\). Suppose \(p^2=r\) and
put \(m=(n-1)/2\). The word commutes with \(r\). Take the half
interval \(I=\{-m,\ldots,m\}\), which \(Hr\) exchanges with its
complement. The preceding argument gives one through edge of \(A\)
and one of \(HAH\). Consequently \(p^{-1}(I)\) is a single cyclic
interval of length \(n\). It contains zero and is invariant under
\(r\). The centers of the two invariant intervals of this length
are zero and \(n\), so \(p(I)=I\). On \(1,\ldots,m\), the induced
permutation \(j\mapsto |p(j)|\), with signed representatives in
\(I\), is a parity-preserving involution. A fixed point would give
\(p(j)=\pm j\) and \(p^2(j)=j\), which is impossible. Therefore the
even and odd subsets both have even cardinality, forcing \(4\mid m\).

The restriction is sharp in every eligible order. A witness at order
nine, read from position zero, is
\[
(0,3,4,17,16,11,10,5,6,9,12,13,8,7,2,1,14,15).
\]
Given a word \(p\) counted by \(Q_n\), choose a
site other than \(0,n\). Its orbit under \(p\) has size four, and
its translate by \(H\) is disjoint because \(H\) changes parity
and \(p\) preserves it. Replace each of these eight crossings by
three consecutive crossings in a small disk. Along the second
curve the three may be read in either increasing or decreasing
order along the first. Both choices are realized by a local
zigzag of one arc across the other. For a block at position \(i\)
in the word, choose the within-block sign \(\epsilon_i\) so that
\(\epsilon_{p(i)}=-\epsilon_i\) and
\(\epsilon_{H(i)}=\epsilon_i\). These prescriptions are consistent
on the two four-cycles. On every enlarged block the square of the
new permutation reverses order. It acts on the blocks as \(r\),
and the half turn still commutes with the word. Thus \(Q_{n+8}>0\).
In particular, for odd \(n>1\),
\[
 Q_n>0\quad\Longleftrightarrow\quad n\equiv1\pmod8.
\]

\subsection*{Computing the fixed sets}
Write \(p(2i)=2\sigma(i)\) for \(i\in\mathbb Z/n\mathbb Z\) and
put \(b=(n-1)/2\). Commutation with \(H\) determines the odd entries.
For \(E_n\), enumerate involutions \(\sigma\) fixing zero and
require the \(n\) chords
\[
  \bigl(2\sigma(i),\,2\sigma(i-b)+n\bmod 2n\bigr)               \tag{A3}
\]
to be noncrossing. These are precisely the upper arcs. Their
antipodal images are all the lower arcs, and their alternating
traversal is the full permutation \(p\). For \(Q_n\), replace the
involution condition by \(\sigma^2(i)=-i\). This pairs the
transpositions of negation into blocks of two, each admitting two
four-cycles. A second computation uses the invariant interval \(I\)
instead. It pairs absolute values of the same parity and assigns
one of the two signs to the image of the first value in each pair.

In the involution recursion each choice fixes one symbol or a
transposition, and a chord can be tested as soon as both images
are known. For larger orders the recursion reuses completions.
The free domain symbols form directed runs under addition by \(b\).
A run of length \(\ell\) has two physical sites per free symbol
and one anchored site at each end. Their order along the run
determines the obligations in (A3), while the physical sites that
remain in each face of the disk inherit a cyclic order. The lengths
of the runs and these face orders specify the completion problem.
Renaming runs and independently reversing face orders preserve the
set of admissible completions up to bijection. Counts are reused
only after comparison of the entire encoded residual data.

One large branch has an independent check in the preceding odd
order. For \(b=(n-1)/2\), the branch \(\sigma(b)=b+1\) contributes
exactly \(E_{n-2}\), with \(E_1=1\). On a word in this branch,
\(p\) interchanges \(-1,1\) and \(n-1,n+1\), and fixes \(0,n\).
Delete the four interchanged sites from both copies of the word.
The resulting permutation remains an involution and commutes with
the new half turn. To see noncrossing directly, delete the two
adjacent arcs \((0,1),(n-1,n)\) of the old lower matching. Give
the surviving sites \(-1,n+1\) the labels of the two new roots.
This is the new upper matching, as is seen by reading successive
entries of the shortened word. The same operation on the old upper
matching, deleting \((-1,0),(n,n+1)\), gives the new lower one.
Inserting the adjacent arcs and restoring the two reversed triples
is the inverse. Thus the identity holds for each order.

Direct generation of all noncrossing matchings through \(n=13\)
checked all the group generators and both fixed-set counts in odd
orders \(3,\ldots,13\). The plain involution search reached order
27; it agreed with the calculation using residual states throughout
the overlap. The two enumerations of the fourth-order symmetry
agreed at \(n=9,17,\ldots,65\), where their common values are
\[
  2,\ 10,\ 70,\ 568,\ 5022,\ 46964,\ 456932,\ 4578628.
\]
The following values illustrate the scale of the self-dual count.
\[
\begin{array}{r|r@{\qquad}r|r}
n&E_n&n&E_n\\ \hline
 3&2&21&2\,694\,816\\
 5&8&23&13\,862\,948\\
 7&36&25&71\,591\,504\\
 9&170&27&370\,926\,926\\
11&826&29&1\,927\,211\,540\\
13&4\,084&31&10\,037\,423\,838\\
15&20\,452&33&52\,388\,267\,330\\
17&103\,374&35&273\,940\,480\,428\\
19&526\,298&37&1\,434\,822\,175\,554
\end{array}
\]
At order 39 the 38 branches of the involution search give
\[
 E_{39}=7\,526\,335\,601\,910.
\]
The direct contribution of branch 20 was \(1\,434\,822\,175\,554\),
equal to the independently computed \(E_{37}\) as predicted by
contraction. Branches 21 through 38 were rerun with different
memoization levels and an open-addressed table. The two computations
agreed on each branch. The full run took 5216 seconds and reached a
peak resident memory of 2.32 GiB.
Since \(D_{20}=566\,068\,240\), formula (A2) predicts
\(P_{39}\equiv-2E_{39}\equiv4\pmod8\). This agrees with the
full projective count. In fact, (A1) was verified in every height
and (A2) against the full totals in all nineteen odd orders from
3 through 39.

\begin{thebibliography}{9}\small
\bibitem{legendre} S.~Legendre, Foldings and meanders, \emph{Australas. J. Combin.} \textbf{58} (2014), 275--291. \href{https://ajc.maths.uq.edu.au/pdf/58/ajc_v58_p275.pdf}{Article}.
\bibitem{folding} J.~Sawada and R.~Li, Stamp foldings, semi-meanders, and open meanders: fast generation algorithms, \emph{Electron. J. Combin.} \textbf{19}(2) (2012), P43. \href{https://doi.org/10.37236/2404}{Publisher}.
\bibitem{how} A.~Howroyd, symmetry formula for closed meanders, \href{https://oeis.org/A077460}{OEIS A077460} (2015).
\bibitem{hogan} H.~Hogan, \href{https://pypi.org/project/mapFolding/0.40.1/}{\emph{mapFolding} 0.40.1}, 17 November 2026.
\bibitem{lacroix} M.~A.~La Croix, \emph{Approaches to the Enumerative Theory of Meanders}, University of Waterloo (2003). \href{https://www.math.uwaterloo.ca/~malacroi/Latex/Meanders.pdf}{Thesis}.
\bibitem{tables} OEIS Foundation, \href{https://oeis.org/A005315}{A005315}, \href{https://oeis.org/A000136}{A000136}, \href{https://oeis.org/A005316}{A005316} and \href{https://oeis.org/A060206}{A060206}.
\end{thebibliography}
\end{document}
